% \begin{savequote}
% \includegraphics[width=5cm]{./images/nasa-pd/154960main_image_feature_638_ys_full-00800.jpg}
% 
% Bild: NASA
% \end{savequote}
\def\ChapterCode{VIT}
%%%%%%%%%%%%%%%%%%%%%%%%%%%%%%%%%%%%%%%%%%%%%%%%%%%%%%%%%%%
%%%%%%%%%%%%%%%%%%%%%%%%%%%%%%%%%%%%%%%%%%%%%%%%%%%%%%%%%%%
%%%%%%%%%%%%%%%%%%%%%%%%%%%%%%%%%%%%%%%%%%%%%%%%%%%%%%%%%%%
%%%%%%%%%%%%%%%%%%%%%%%%%%%%%%%%%%%%%%%%%%%%%%%%%%%%%%%%%%%
%%%%%%%%%%%%%%%%%%%%%%%%%%%%%%%%%%%%%%%%%%%%%%%%%%%%%%%%%%%
\chapter
[\CO{[VIT]} Vitalfunktionen]
{Vitalfunktionen}

% \AasRsN{RS~V}


% \emph{Maintainer: Sebastian Gabriel}

% \AuthNotes{
% Changelog:
% 
% \begin{description}
%     \item[2009-02-25] Start Changelog. Kapitel mit Korr
%     Koch korr. 
%     \item[2009-mm-dd] in \LaTeX übernommen
%     \item[2009-04-19] anhand der OpenOffice-Version überarbeitet.
% Bilder fehlen noch. 
% \item[2009-05-06] GAB: Abschnitte \enquote{Blutdruck} und
% \enquote{Atemvolumina} aus Anatomie-Teil übernommen.
% \end{description}
% 
% }

\ChapterIntro
{Diese Kapitel behandelt die Vitalfunktionen
des menschlichen Körpers. Sie ermöglichen die Funktion des
Körpers. Es gibt grundlegende Vitalfunktionen 1. Ordnung
(Bewusstsein, Atmung, Kreislauf), deren Ausfall binnen
kurzem zum Tod führen können, und Vitalfunktionen 2.
Ordnung, deren
Störung längere Zeit toleriert werden kann.

Die Beurteilung und Erhaltung von Vitalfunktionen ist eine
wesentliche Aufgabe in der Sanitätshilfe.
}{%
\begin{description}
  \item[Maintainer:] S. Gabriel
  \item[Autoren:] Diverse
  \item[Reviewer:] Standard-Reviewprozess
  \item[Version:] {\VarDocVersionTypeName} ({\VarDocVersionTypeDescription})
  \item[SHA1:] {\DocGitCommit}
\end{description}

{\TxTeilkapitelAASS}
}

\ChapterEnd